\chapter{现代循环神经网络}

门控让隐状态学会何时写入、何时遗忘。GRU 与 LSTM 用可学习的门缓解长程梯度问题，
再堆成深层或双向结构。编码器--解码器把变长输入映成变长输出，束搜索在词表上近似最大似然序列。

\section{门控循环单元（GRU）}

普通循环中 $\prod_j\partial h_j/\partial h_{j-1}$ 容易消失或爆炸。
若早期词元必须影响很晚的预测，就需要能够长期保存 $h$ 的机制；
若中间词元无关，则应学会跳过更新。GRU 用重置门与更新门实现这两种行为。

\subsection{门控隐状态}

门的取值在 $(0,1)$ 中，由当前输入与上一隐状态共同决定，因此何时重置、何时拷贝都可以学习。
每个门分量表示对该维摘要“放行多少”：靠近 $1$ 几乎全放，靠近 $0$ 几乎全挡。

\subsubsection{重置门和更新门}

对批量输入 $\bm{X}_t$ 与上一隐状态 $\bm{H}_{t-1}$，
\begin{align*}
\bm{R}_t
&=
\sigma(\bm{X}_t\bm{W}_{xr}+\bm{H}_{t-1}\bm{W}_{hr}+\bm{b}_r),
\\
\bm{Z}_t
&=
\sigma(\bm{X}_t\bm{W}_{xz}+\bm{H}_{t-1}\bm{W}_{hz}+\bm{b}_z).
\end{align*}
sigmoid 把线性打分压到 $(0,1)$。
$\bm{R}_t$ 的每个分量是重置比例：靠近 $0$ 表示几乎忘掉上一摘要的该维，靠近 $1$ 表示几乎原样用于计算候选。
$\bm{Z}_t$ 的每个分量是更新比例：靠近 $1$ 表示本维倾向原样保留 $\bm{H}_{t-1}$，靠近 $0$ 表示倾向换成新候选。
二者都不是隐状态本身，而是 $[0,1]$ 上的开关程度。

\subsubsection{候选隐状态}

重置后的历史再与当前输入混合：
\[
\tilde{\bm{H}}_t
=
\tanh\bigl(\bm{X}_t\bm{W}_{xh}+(\bm{R}_t\odot\bm{H}_{t-1})\bm{W}_{hh}+\bm{b}_h\bigr).
\]
$\bm{R}_t\odot\bm{H}_{t-1}$ 是按维缩放后的旧摘要；$\tilde{\bm{H}}_t$ 的分量在 $(-1,1)$，是“若完全改写，新摘要该写成什么”。
$\bm{R}_t=\bm{1}$ 时退化为普通候选隐状态；$\bm{R}_t=\bm{0}$ 时候选只看 $\bm{X}_t$，旧历史不进入候选。

\subsubsection{隐状态}

更新门在旧状态与候选之间插值：
\[
\bm{H}_t
=
\bm{Z}_t\odot\bm{H}_{t-1}+(1-\bm{Z}_t)\odot\tilde{\bm{H}}_t.
\]
$\bm{H}_t$ 仍是到时刻 $t$ 为止的序列压缩摘要，与 GRU 的 $h$ 同形状。
每个分量是旧摘要与新候选的凸组合：$\bm{Z}_t$ 靠近 $1$ 则几乎拷贝旧摘要，靠近 $0$ 则几乎采用候选。
$\bm{Z}_t=\bm{1}$ 时 $\bm{H}_t=\bm{H}_{t-1}$，梯度可沿时间原路回传；
$\bm{Z}_t=\bm{0}$ 时完全改写。

\subsection{从零开始实现}

门控递推与上一章字符语言模型相同，只是 $h_t$ 的更新换成上面三式，输出层仍是
$\bm{O}_t=\bm{H}_t\bm{W}_{hq}+\bm{b}_q$。
$\bm{O}_t$ 的分量是词表上的 logits；softmax 后才是下一词概率。

\subsubsection{初始化模型参数}

高斯初始化权重、零偏置。需要实例化
$(\bm{W}_{xr},\bm{W}_{hr},\bm{b}_r)$、
$(\bm{W}_{xz},\bm{W}_{hz},\bm{b}_z)$、
$(\bm{W}_{xh},\bm{W}_{hh},\bm{b}_h)$
以及输出层 $(\bm{W}_{hq},\bm{b}_q)$。隐藏宽度仍由 $h$ 给出，$h$ 是摘要向量的维数。

\subsubsection{定义模型}

$\bm{H}_0=\bm{0}$，表示空摘要。逐步计算 $(\bm{R}_t,\bm{Z}_t,\tilde{\bm{H}}_t,\bm{H}_t,\bm{O}_t)$，
参数在所有 $t$ 上共享。映射与定义式逐项相同：门在 $(0,1)$，摘要在 $(-1,1)$ 量级，logits 无界。

\subsubsection{训练与预测}

损失、裁剪、困惑度与前缀续写均沿用循环神经网络从零实现中的流程。
训练后比较训练集困惑度（有效分支数）以及给定前缀时的续写序列。

\subsection{简洁实现}

框架中的 GRU 层实现同一组门。前向仍是
\[
(\bm{H}_{1:T},\bm{H}_T)=\operatorname{GRU}(\bm{X}_{1:T},\bm{H}_0),
\]
$\bm{H}_{1:T}$ 的每个切片是对应时刻的压缩摘要，$\bm{H}_T$ 是读完整段后的最终摘要。
再接输出层；数值上应与从零实现一致，只是算子被编译。

\subsection{小结}

GRU 的状态更新
\[
\bm{H}_t=\bm{Z}_t\odot\bm{H}_{t-1}+(1-\bm{Z}_t)\odot\tilde{\bm{H}}_t
\]
用 $\bm{Z}_t\in(0,1)$ 决定保留多少旧摘要、用 $\bm{R}_t\in(0,1)$ 决定候选里用多少旧历史；
$\bm{H}_t$ 仍是迄今序列的压缩摘要。两门全开或全关时分别覆盖拷贝与普通 RNN。

\subsection{练习}

核验仅用 $x_{t'}$ 预测更晚输出时 $\bm{R},\bm{Z}$ 的理想取值，
以及只保留一扇门时递推如何退化。
长期记忆应对应 $\bm{Z}\approx\bm{1}$（少改写摘要），忽略一段历史应对应 $\bm{R}\approx\bm{0}$。

\section{长短期记忆网络（LSTM）}

LSTM 在隐状态之外再设记忆元 $\bm{C}_t$，用三扇门分别控制写入、遗忘与读出。

\subsection{门控记忆元}

记忆元与 $\bm{H}_t$ 同形，但不直接进入输出层。输入门决定吸收候选，
遗忘门决定丢弃旧记忆，输出门决定暴露多少记忆给 $h_t$。
三扇门的每个分量都在 $(0,1)$：数值就是“做该操作的比例”。

\subsubsection{输入门、忘记门和输出门}

\begin{align*}
\bm{I}_t
&=
\sigma(\bm{X}_t\bm{W}_{xi}+\bm{H}_{t-1}\bm{W}_{hi}+\bm{b}_i),
\\
\bm{F}_t
&=
\sigma(\bm{X}_t\bm{W}_{xf}+\bm{H}_{t-1}\bm{W}_{hf}+\bm{b}_f),
\\
\bm{O}_t
&=
\sigma(\bm{X}_t\bm{W}_{xo}+\bm{H}_{t-1}\bm{W}_{ho}+\bm{b}_o).
\end{align*}
三者均在 $(0,1)$ 中逐元素取值。
$I_{t,j}$ 是第 $j$ 维写入多少新候选：$1$ 表示全部写入，$0$ 表示不写入。
$F_{t,j}$ 是第 $j$ 维遗忘多少旧记忆：$1$ 表示全部保留旧记忆，$0$ 表示全部忘掉。
注意：遗忘门靠近 $1$ 是“少忘”，靠近 $0$ 是“多忘”。
$O_{t,j}$ 是第 $j$ 维向外输出多少记忆：$1$ 表示该维记忆充分暴露给隐状态，$0$ 表示对外隐藏。

\subsubsection{候选记忆元}

\[
\tilde{\bm{C}}_t
=
\tanh(\bm{X}_t\bm{W}_{xc}+\bm{H}_{t-1}\bm{W}_{hc}+\bm{b}_c).
\]
$\tanh$ 把候选限制在 $(-1,1)$。每个分量是“若输入门打开，打算写入记忆的新内容”，不是门。

\subsubsection{记忆元}

\[
\bm{C}_t
=
\bm{F}_t\odot\bm{C}_{t-1}+\bm{I}_t\odot\tilde{\bm{C}}_t.
\]
$\bm{C}_t$ 是内部记忆：旧记忆按遗忘门缩放，新候选按输入门缩放，再相加。
每个分量仍是记忆内容（可正可负），不是 $(0,1)$ 的门。
$\bm{F}_t=\bm{1},\bm{I}_t=\bm{0}$ 时记忆原样传递，梯度可沿 $\bm{C}$ 传播。

\subsubsection{隐状态}

\[
\bm{H}_t=\bm{O}_t\odot\tanh(\bm{C}_t).
\]
先把记忆压到 $(-1,1)$，再按输出门决定对外暴露多少。
$\bm{H}_t$ 是对外可见的序列压缩摘要，供输出层使用；
$\bm{C}_t$ 是内部状态，不直接解码。
只有 $\bm{H}_t$ 进入 $\bm{O}_t^{\text{out}}=\bm{H}_t\bm{W}_{hq}+\bm{b}_q$，后者是词表 logits。

\subsection{从零开始实现}

隐状态现在是二元组 $(\bm{H}_t,\bm{C}_t)$，其余训练回路不变。
$\bm{H}_t$ 是摘要，$\bm{C}_t$ 是记忆。

\subsubsection{初始化模型参数}

除三扇门与候选记忆的权重偏置外，输出层维度仍等于词表大小。
权重取标准差 $0.01$ 的高斯，偏置为零。
$0.01$ 是初始权重的标准差，使初始打分接近 $0$。

\subsubsection{定义模型}

初始化 $(\bm{H}_0,\bm{C}_0)=(\bm{0},\bm{0})$，表示空摘要与空记忆。逐步计算
$(\bm{I}_t,\bm{F}_t,\bm{O}_t,\tilde{\bm{C}}_t,\bm{C}_t,\bm{H}_t)$，
仅把 $\bm{H}_t$ 映到词表 logits。
门在 $(0,1)$，候选与 $\tanh(\bm{C}_t)$ 在 $(-1,1)$。

\subsubsection{训练和预测}

与 GRU 相同的困惑度目标与前缀生成。记忆元在预测时随 $\bm{H}$ 一起递推，但不直接解码。
困惑度仍是每个位置有效分支数的几何平均。

\subsection{简洁实现}

\[
\bigl((\bm{H}_{1:T},\bm{C}_T),\bm{H}_T\bigr)
=
\operatorname{LSTM}(\bm{X}_{1:T},(\bm{H}_0,\bm{C}_0)).
\]
返回的 $\bm{H}_{1:T}$ 是各时刻对外摘要，$\bm{C}_T$ 是最终内部记忆。
这是带状态控制的隐变量自回归模型；长程依赖仍使逐步展开代价随 $T$ 增长，$T$ 是时间步数。

\subsection{小结}

LSTM 用
\[
\bm{C}_t=\bm{F}_t\odot\bm{C}_{t-1}+\bm{I}_t\odot\tilde{\bm{C}}_t,
\qquad
\bm{H}_t=\bm{O}_t\odot\tanh(\bm{C}_t)
\]
把长期记忆与对外隐状态分开。
$\bm{F}_t,\bm{I}_t,\bm{O}_t\in(0,1)$ 分别是遗忘、写入、读出的比例；
$\bm{H}_t$ 是迄今序列的压缩摘要。从而缓解梯度消失。

\subsection{练习}

核验 $\tilde{\bm{C}}_t$ 已经过 $\tanh$ 后输出为何再套 $\tanh$，
以及在相同 $h$ 下 GRU、LSTM 与普通 RNN 的逐步乘法次数。
第二次 $\tanh$ 把记忆幅度限制在 $(-1,1)$ 再交给输出门；$h$ 是摘要维数。

\section{深度循环神经网络}

单层循环的 $\phi$ 可能不够灵活。把若干循环层堆叠，使不同层刻画不同时间尺度。

\subsection{函数依赖关系}

令 $\bm{H}_t^{(0)}=\bm{X}_t$。第 $\ell=1,\ldots,L$ 层
\[
\bm{H}_t^{(\ell)}
=
\phi_\ell\bigl(
\bm{H}_t^{(\ell-1)}\bm{W}_{xh}^{(\ell)}
+
\bm{H}_{t-1}^{(\ell)}\bm{W}_{hh}^{(\ell)}
+
\bm{b}_h^{(\ell)}
\bigr).
\]
$\bm{H}_t^{(\ell)}$ 是第 $\ell$ 层在时刻 $t$ 的压缩摘要：它混合了下层同一时刻的摘要与本层上一时刻的摘要。
$L$ 是层数。
顶层再接到输出
\[
\bm{O}_t=\bm{H}_t^{(L)}\bm{W}_{hq}+\bm{b}_q.
\]
$\bm{O}_t$ 是词表打分。
于是信息同时沿时间传到 $\bm{H}_{t+1}^{(\ell)}$，并沿深度传到 $\bm{H}_t^{(\ell+1)}$。

\subsection{简洁实现}

以 LSTM 为例，仅增加层数 $L$。词表大小决定输入输出维，隐藏宽度仍取 $h$，
层间按上式串联。GRU 或普通 RNN 只需更换 $\phi_\ell$ 所对应的单元。
每层仍有自己的门（若用 LSTM），门值含义不变：$(0,1)$ 上的遗忘/写入/读出比例。

\subsection{训练与预测}

多层逐步展开使每步代价约为单层的 $L$ 倍，收敛对学习率与裁剪更敏感。
预测时逐层递推全部 $\bm{H}_t^{(\ell)}$，每一层都维护到当前为止的摘要。

\subsection{小结}

深度循环的依赖是
\[
\bm{H}_t^{(\ell)}=
\phi_\ell(\bm{H}_t^{(\ell-1)}\bm{W}_{xh}^{(\ell)}+\bm{H}_{t-1}^{(\ell)}\bm{W}_{hh}^{(\ell)}+\bm{b}_h^{(\ell)}),
\]
$\bm{H}_t^{(\ell)}$ 是第 $\ell$ 层对迄今历史的压缩摘要。
同一套门控单元可以替换 $\phi_\ell$。

\subsection{练习}

核验两层从零展开时 $\bm{H}_t^{(1)}$ 如何作为 $\bm{H}_t^{(2)}$ 的输入，
以及用 GRU 替换 LSTM 后困惑度与逐步代价如何变化。
下层摘要是上层的输入特征，不是词表概率。

\section{双向循环神经网络}

单向循环只能看过去。填空、标注等任务需要未来上下文，因此再设一条反向循环。

\subsection{隐马尔可夫模型中的动态规划}

隐状态与观测的联合分布为
\[
P(x_1,\ldots,x_T,h_1,\ldots,h_T)
=
\prod_{t=1}^{T}P(h_t \mid h_{t-1})P(x_t \mid h_t),
\]
其中 $P(h_1 \mid h_0)=P(h_1)$。
每个因子是转移概率或发射概率，均在 $[0,1]$；乘积是隐链与观测同时出现的联合概率。
边际化全部 $h$ 时，前向量
\[
\pi_{t+1}(h_{t+1})
=
\sum_{h_t}\pi_t(h_t)P(x_t \mid h_t)P(h_{t+1} \mid h_t)
\]
把过去的贡献压缩成关于 $h_{t+1}$ 的函数：$\pi_{t+1}(h)$ 是“到 $t+1$ 为止、隐状态为 $h$”的非归一化质量，非负。
最终
\[
P(x_1,\ldots,x_T)=\sum_{h_T}\pi_T(h_T)P(x_T \mid h_T).
\]
该和是观测序列的边际似然，在 $[0,1]$。
后向量
\[
\rho_{t-1}(h_{t-1})
=
\sum_{h_t}P(h_t \mid h_{t-1})P(x_t \mid h_t)\rho_t(h_t)
\]
压缩未来：$\rho_{t-1}(h)$ 是从 $h_{t-1}=h$ 出发解释后续观测的质量。
两者合起来给出缺省位置的条件
\[
P(x_j \mid x_{-j})
\propto
\sum_{h_j}\pi_j(h_j)\rho_j(h_j)P(x_j \mid h_j).
\]
归一化后是“已知其余观测时，位置 $j$ 取各词元”的概率，在 $[0,1]$ 且对 $x_j$ 求和为 $1$。

\subsection{双向模型}

用一条前向 RNN 扮演 $\pi$，一条后向 RNN 扮演 $\rho$，在每个 $t$ 拼接两边的隐状态。

\subsubsection{定义}

\begin{align*}
\overrightarrow{\bm{H}}_t
&=
\phi(\bm{X}_t\bm{W}_{xh}^{(f)}+\overrightarrow{\bm{H}}_{t-1}\bm{W}_{hh}^{(f)}+\bm{b}_h^{(f)}),
\\
\overleftarrow{\bm{H}}_t
&=
\phi(\bm{X}_t\bm{W}_{xh}^{(b)}+\overleftarrow{\bm{H}}_{t+1}\bm{W}_{hh}^{(b)}+\bm{b}_h^{(b)}).
\end{align*}
$\overrightarrow{\bm{H}}_t$ 是只看 $x_1,\ldots,x_t$ 的压缩摘要；
$\overleftarrow{\bm{H}}_t$ 是只看 $x_t,\ldots,x_T$ 的压缩摘要。
拼接 $\bm{H}_t=[\overrightarrow{\bm{H}}_t;\overleftarrow{\bm{H}}_t]$ 后
\[
\bm{O}_t=\bm{H}_t\bm{W}_{hq}+\bm{b}_q.
\]
$\bm{H}_t$ 是双向上下文的联合摘要；$\bm{O}_t$ 是该位置的标签打分。

\subsubsection{模型的计算代价及其应用}

每个 $t$ 的输出同时依赖 $x_1,\ldots,x_T$，因此不能用于“只根据过去预测下一个词元”。
前向需要跑完反向递归，反向传播的链大约加倍，训练代价高。适用场合是编码整句或双向上下文标注。

\subsection{双向循环神经网络的错误应用}

若把双向模型当作语言模型训练，训练时未来词元可见，测试时不可见，
条件分布 $P(x_t \mid x_{<t})$ 被错误地换成 $P(x_t \mid x_{\neq t})$，续写会崩溃。
前者是因果下一词概率，后者是填空概率，二者都在 $[0,1]$ 但条件不同。

\subsection{小结}

双向更新
\[
\bm{O}_t=\bigl[\overrightarrow{\bm{H}}_t;\overleftarrow{\bm{H}}_t\bigr]\bm{W}_{hq}+\bm{b}_q
\]
对应 HMM 中的前向--后向量；左右两段都是压缩摘要，合起来编码双向上下文，而不能当作因果语言模型。

\subsection{练习}

核验两个方向隐藏宽度不同时 $\bm{H}_t$ 的拼接维，
以及如何用双向上下文表示多义词在句中的向量。
拼接维是两方向摘要维数之和，该整数是向量长度，不是概率。

\section{机器翻译与数据集}

机器翻译把源语言序列映到目标语言序列，是变长到变长的序列转换。
后续编码器--解码器都在这对平行语料上训练。

\subsection{下载和预处理数据集}

英--法句对经制表符分隔。预处理把不间断空格换成普通空格、字母小写，
并在单词与标点之间插入空格，得到可切分的字符串对 $(s,t)$。
此时仍是文本，不是张量。

\subsection{词元化}

与字符级语言模型不同，这里按单词（及标点）切分。第 $i$ 对变成
\[
\operatorname{source}^{(i)}=(x_1,\ldots,x_{T}),
\qquad
\operatorname{target}^{(i)}=(y_1,\ldots,y_{T'}).
\]
每个 $x,y$ 稍后会被映成词表下标；$T,T'$ 是源、目标词元个数。
多数句对的词元数小于 $20$，$20$ 是长度阈值。

\subsection{词表}

源、目标各建一个词表。单词级使 $|\mathcal{V}|$ 远大于字符级，
故把频次小于 $2$ 的项并入 $\texttt{<unk>}$，并加入
$\texttt{<pad>}$、$\texttt{<bos>}$、$\texttt{<eos>}$。
$2$ 是最低频次；$|\mathcal{V}|$ 是词表大小，决定 softmax 有几类。

\subsection{加载数据集}

句长不同。截断与填充把每条序列补/切到固定 $n_{\text{steps}}$，
同时记录有效长度 $T_{\text{valid}}\le n_{\text{steps}}$，
以便后面的损失与注意力忽略填充。
$n_{\text{steps}}$ 与 $T_{\text{valid}}$ 都是词元个数（轴长度），填充位置上的下标是 $\texttt{<pad>}$ 的编号，不是语义词。

\subsection{训练模型}

数据迭代器每次返回一个小批量的源下标、源有效长度、目标下标与目标有效长度，
以及两个词表。这给出监督学习对 $(\bm{X},\bm{Y})$：
每个条目是词元编号，有效长度是有几个非填充位置。

\subsection{小结}

平行语料经过词元化与填充后，每个样本是长度同为 $n_{\text{steps}}$ 的整数对
$(\bm{x},\bm{y})$，低频率词并入未知词元以控制 $|\mathcal{V}|$。
这些整数是类别下标，不是嵌入坐标。

\subsection{练习}

核验句对数如何改变两个词表大小，以及无空格语言上单词切分是否仍有良定义的 $\mathcal{Z}$。
句对数是样本个数；$|\mathcal{V}|$ 随词种增加。

\section{编码器-解码器架构}

输入长度 $T$ 与输出长度 $T'$ 都可以变。编码器把变长输入压成固定形状的状态，
解码器再把该状态展开成变长输出。

\subsection{编码器}

编码器是映射
\[
\operatorname{Enc}:\ \mathcal{X}^{T}\to\mathcal{C},
\qquad
\bm{X}\mapsto\bm{c}.
\]
$\bm{X}$ 是长度为 $T$ 的源词元（或嵌入）序列；$\bm{c}$ 是固定形状的上下文向量（或张量），
不随 $T$ 改变轴长度，含义是整句源序列的压缩摘要。
任何具体网络只要实现这一接口即可。

\subsection{解码器}

解码器在每步消费已经生成的词元与编码状态：
\[
\operatorname{Dec}:\ \mathcal{C}\times\mathcal{Y}^{t'-1}\to\mathcal{Y},
\qquad
(\bm{c},y_{1:t'-1})\mapsto y_{t'}.
\]
输出 $y_{t'}$ 是目标词表上的下一个词元（或先给出该词的概率再取样）。
$\operatorname{init\_state}$ 把编码器输出转换成解码器的初始状态，必要时使用源序列有效长度。

\subsection{合并编码器和解码器}

前向先算 $\bm{c}=\operatorname{Enc}(\bm{X})$，再逐步
\[
y_{t'}=\operatorname{Dec}(\bm{c},y_{1:t'-1}),
\qquad t'=1,\ldots,T'.
\]
每步 $y_{t'}$ 是一个目标词元；$T'$ 是目标长度。
具有隐状态的网络自然适合携带 $\bm{c}$。下一节用两个循环网络实现该架构。

\subsection{小结}

编码器--解码器把变长到变长写成
\[
\bm{c}=\operatorname{Enc}(\bm{X}),
\qquad
y_{1:T'}=\operatorname{Dec}(\bm{c}),
\]
$\bm{c}$ 是源句摘要，$y_{1:T'}$ 是目标词元序列。
因此适用于机器翻译这类序列转换。

\subsection{练习}

核验 $\operatorname{Enc}$ 与 $\operatorname{Dec}$ 不必同一类网络，
以及哪些其他变长到变长任务可用同一接口。
接口约束的是输入输出对象，不是内部用卷积还是循环。

\section{序列到序列学习（seq2seq）}

用循环神经网络实现上一节的两个映射：编码器消费 $\bm{x}_{1:T}$，
解码器生成 $\bm{y}_{1:T'}$。

\subsection{编码器}

逐步
\[
\bm{h}_t=f(\bm{x}_t,\bm{h}_{t-1}).
\]
$\bm{h}_t$ 是源序列读到第 $t$ 个词元时的压缩摘要。
上下文变量把全部隐状态汇总为固定向量
\[
\bm{c}=q(\bm{h}_1,\ldots,\bm{h}_T).
\]
$\bm{c}$ 是整句源端摘要。最常见的取法是 $q$ 只保留 $\bm{h}_T$，或取多层最终隐状态的拼接；
拼接增加的是向量长度，不是把摘要数值加总。

\subsection{解码器}

解码器隐状态
\[
\bm{s}_{t'}=g(y_{t'-1},\bm{c},\bm{s}_{t'-1}),
\]
$\bm{s}_{t'}$ 是目标端到 $t'$ 为止的压缩摘要，同时注入源摘要 $\bm{c}$。
再由 $\bm{s}_{t'}$ 产生 $P(y_{t'}\mid y_{1:t'-1},\bm{c})$：
词表上的条件概率，在 $[0,1]$ 且和为 $1$。
训练时 $y_{t'-1}$ 来自真实标签，预测时来自上一时刻的输出。

\subsection{损失函数}

每步在词表上做 softmax 交叉熵。填充位置不应计入损失。
有效长度掩码把越界项置零，等价于
\[
L
=
\frac{1}{\sum_{t'}\mathbb{I}\{t'\le T'_{\text{valid}}\}}
\sum_{t'=1}^{n_{\text{steps}}}
\mathbb{I}\{t'\le T'_{\text{valid}}\}\,
\ell(y_{t'},\hat{y}_{t'}).
\]
指示函数取值 $\{0,1\}$：$1$ 表示该步是真实词元。
$L$ 是有效位置上交叉熵的平均，标量，单位与逐步负对数概率相同；越小表示给正确目标词的概率越高。

\subsection{训练}

解码器输入是 $\texttt{<bos>}$ 与去掉 $\texttt{<eos>}$ 的真实目标序列拼接，
即强制教学。也可以改用模型自己的 $\hat{y}_{t'-1}$。编码器可多层，解码器初始 $\bm{s}_0$ 由 $\bm{c}$ 给出，
即用源摘要初始化目标端摘要。

\subsection{预测}

从 $\texttt{<bos>}$ 开始，每步把 $\hat{y}_{t'}$ 送回解码器，直到输出 $\texttt{<eos>}$
或达到最大长度。这是贪心生成；下一节讨论束搜索。
每步选出的是词元编号，其依据是 $P(y_{t'}\mid\cdot)$ 最大。

\subsection{预测序列的评估}

BLEU 比较预测与标签的 $n$ 元匹配率 $p_n$，并惩罚过短输出：
\[
\exp\Biggl(
\min\Bigl(0,1-\frac{\mathrm{len}_{\mathrm{label}}}{\mathrm{len}_{\mathrm{pred}}}\Bigr)
\Biggr)
\prod_{n=1}^{k}p_n^{1/2^{n}}.
\]
$p_n\in[0,1]$ 是长度为 $n$ 的片段命中率；长度比无量纲。
整个 BLEU 在 $[0,1]$ 附近（指数惩罚项 $\le 1$），越接近 $1$ 表示 $n$ 元重叠越好且长度不过短。

\subsection{小结}

seq2seq 用
\[
\bm{c}=q(\bm{h}_{1:T}),
\qquad
\bm{s}_{t'}=g(y_{t'-1},\bm{c},\bm{s}_{t'-1})
\]
实现编码器--解码器：$\bm{c}$ 与 $\bm{s}_{t'}$ 分别是源、目标侧的压缩摘要。
训练用强制教学，评价用 BLEU。

\subsection{练习}

核验去掉掩码后填充位置如何污染 $L$，
以及编码器与解码器隐藏宽度不同时 $\bm{s}_0$ 应如何由 $\bm{c}$ 初始化。
填充位置的交叉熵不对应真实词，会把平均损失这一标量拉偏。

\section{束搜索}

解码要在词表 $\mathcal{Y}$ 上找一条序列最大化
$P(y_1,\ldots,y_{T'}\mid\bm{c})$。
该目标是整条目标序列在给定源摘要下的联合概率，在 $[0,1]$。
逐步取最大是贪心；枚举全部序列是穷举；
束搜索在二者之间保留 $k$ 条候选。

\subsection{贪心搜索}

每步
\[
y_{t'}=\operatorname*{argmax}_{y\in\mathcal{Y}}P(y \mid y_1,\ldots,y_{t'-1},\bm{c}).
\]
右边每个 $P$ 是下一词条件概率；argmax 返回概率最大的那个词元，不是把概率当作词。
计算量 $\mathcal{O}(|\mathcal{Y}|T')$，但不能改正早期错误，整体不必最优。

\subsection{穷举搜索}

枚举 $\mathcal{Y}^{T'}$ 中全部序列并取联合概率最大者。
计算量 $\mathcal{O}(|\mathcal{Y}|^{T'})$。例如 $|\mathcal{Y}|=10^{4}$、$T'=10$ 时约 $10^{40}$ 条，不可行。
$10^{40}$ 是候选序列条数。

\subsection{束搜索}

束宽为 $k$ 时，每步只保留条件概率最高的 $k$ 条前缀。$k$ 是条数。例如 $k=2$ 时，
若第一步留下 $A,C$，则第二步比较
\begin{align*}
P(A,y_2 \mid \bm{c})&=P(A \mid \bm{c})P(y_2 \mid A,\bm{c}),
\\
P(C,y_2 \mid \bm{c})&=P(C \mid \bm{c})P(y_2 \mid C,\bm{c}),
\end{align*}
再留下两条。左边是长度为 $2$ 的前缀的联合概率，等于逐步条件概率之积，仍在 $[0,1]$。
第三步同理，
\begin{align*}
P(A,B,y_3 \mid \bm{c})&=P(A,B \mid \bm{c})P(y_3 \mid A,B,\bm{c}),
\\
P(C,E,y_3 \mid \bm{c})&=P(C,E \mid \bm{c})P(y_3 \mid C,E,\bm{c}).
\end{align*}
最终在候选完整序列上用长度归一化分数
\[
\frac{1}{L^{\alpha}}\log P(y_1,\ldots,y_L \mid \bm{c})
=
\frac{1}{L^{\alpha}}\sum_{t'=1}^{L}\log P(y_{t'} \mid y_{1:t'-1},\bm{c})
\]
选出输出。
$\log P(y_1,\ldots,y_L\mid\bm{c})$ 是该词元序列的对数概率（nat 或 bit，随对数底而定），通常为负，越大（越接近 $0$）表示整句越可能。
右边把逐步 $\log P(y_{t'}\mid\cdot)$ 相加，每项是该步真实（或候选）词的对数条件概率。
$L^{\alpha}$ 用长度惩罚过短句，$\alpha$ 无量纲。
$k=1$ 退化为贪心；$k=|\mathcal{Y}|$ 恢复穷举。

\subsection{小结}

束搜索用宽度 $k$ 在
\[
\max_{y_{1:L}}P(y_{1:L}\mid\bm{c})
\]
的精度与 $\mathcal{O}(k|\mathcal{Y}|T')$ 的代价之间折中。
比较时所依据的束分数是词元序列的（长度归一化）对数概率，不是单个词的下标。

\subsection{练习}

核验穷举是否为 $k=|\mathcal{Y}|$ 的束搜索，以及 $k$ 如何改变机器翻译的速度与 BLEU。
$k$ 增大则每步保留的前缀更多，搜索更接近最大对数概率，计算按 $k$ 近似线性增加。
